\subsection{Homomorphisms}

Eerst zullen we aantonen dat \(\varphi(1) = 1\):

Voor een homomorfisme \(\varphi: G_1 \to G_2\) geldt \(\varphi(gh) = \varphi(g)\varphi(h)\).

Nu is \(\varphi(1) = \varphi(1 \cdot 1) = \varphi(1)\varphi(1)\). Vermenigvuldig links met \(\varphi(1)^{-1}\):

\(\varphi(1)^{-1}\varphi(1) = \varphi(1)^{-1}\varphi(1)\varphi(1)\)

\(\implies \varphi(1) = 1\)
\qed

\subsubsection{Proposition 2.1}

\begin{proposition}
    The image of a group homomorphism $\varphi: G_1 \to G_2$ is a subgroup of its codomain $G_2$.

    We have \(\varphi(G_1) \leq G_2\).
\end{proposition}

\begin{proof}
    We tonen aan dat $\varphi(G_1) \leq G_2$ via het deelgroepcriterium: identiteit, geslotenheid, en inverses.

    \textbf{Identiteit:} Er geldt
    \[
        \varphi(e_1)\varphi(e_1) = \varphi(e_1 e_1) = \varphi(e_1).
    \]
    Vermenigvuldig links met $\varphi(e_1)^{-1} \in G_2$:
    \[
        \varphi(e_1) = e_2.
    \]
    Dus $e_2 \in \varphi(G_1)$.

    \textbf{Geslotenheid:} Ask a toddler on the street. (Dit is gewoon toepassen van de definitie van een homomorfisme.)
    
    Neem $a, b \in \varphi(G_1)$, dan bestaan er $g, h \in G_1$ met $a = \varphi(g)$ en $b = \varphi(h)$. Dan
    \[
        ab = \varphi(g)\varphi(h) = \varphi(gh) \in \varphi(G_1),
    \]
    want $gh \in G_1$.

    \textbf{Inverses:} Neem $a \in \varphi(G_1)$, dan bestaat er $g \in G_1$ met $a = \varphi(g)$. Omdat $g^{-1} \in G_1$, geldt
    \[
        \varphi(g)\varphi(g^{-1}) = \varphi(gg^{-1}) = \varphi(e_1) = e_2,
    \]
    dus $\varphi(g^{-1}) = \varphi(g)^{-1} = a^{-1}$. Bijgevolg
    \[
        a^{-1} = \varphi(g^{-1}) \in \varphi(G_1).
    \]

    Aangezien $\varphi(G_1)$ het neutraal element bevat, gesloten is onder de bewerking, en gesloten is onder inverteren, volgt $\varphi(G_1) \leq G_2$.
\end{proof}